\setcounter{equation}{0}

\chapter{Introduction}

\section{Problem Statement}
\lipsum[1]

\section{Problem Description}
\lipsum[1]

\section{Significance of the Problem}
Document clustering seeks to identify groups of documents, from a given text document collection, such that documents in a group are more similar in their theme and documents from different groups are more dissimilar in their meaning.

Due to rapid growth in the usage of internet \cite{warf2020}, the volume of online text without associated class labels; in the form web documents, judicial records, medical health records, operation manuals, research articles, news articles and textbooks is increasing \cite{zhai2016},\cite{bharti2014},\cite{feldman2007},\cite{forman2003}. Text clustering is essential to organize, summarize and to retrieve relevant information automatically from the online text corpora \cite{garcia2020}. 

\section{Challenges in Solving the Problem}
 
Text documents are in unstructured format and are difficult to analyse in their original form. They are to be represented in a structured format for the sake of analysis. Most popular document representation technique is Vector Space Mode(VSM) \cite{zapata2000}. 

There exist complex relationships among terms \cite{zapata2000}. High dimensionality and sparseness of DTM not only increases the time complexity of clustering algorithms but also reduces the clustering performance of the algorithms. 


\section{Assumptions and Limitations of the Work}
The assumptions and limitations of this work are 
\begin{inparaenum}[\itshape i\upshape)]
\item The text documents in the collections are written in English language. This is a reasonable assumption as most of the online text available is in English.
\item Each text document describes a single theme. The documents are assumed to belong to only one category.
\item Each document has reasonable length. The number of terms in a document is 10 \% more than the number of unique terms in the document.
\item Each document has reasonable number of unique terms. If the number of unique terms in a document is less than one percentile of unique terms collection of all documents in a text corpus, the document is not considered for clustering. 
\end{inparaenum}


\section{Organization of the Thesis}
Chapter 2 presents a detailed review of existing solutions for the problem reported in the literature.  Chapter 3 presents details about the proposed solution. Chapter 4 presents the details about the methodology that is followed to do this work. Chapter
5 presents the details about results obtained and analysis of the results. Chapter 6 provides a
detailed conclusion of the present work with a description of the scope for future work.
